%%%%%%%%%%%%%%%%%%%%%
% Documento maestro %
%%%%%%%%%%%%%%%%%%%%%
\documentclass{fime_eng}

%%%%%%%%%%%%%%%%%%%%%%%%%%%%%%%%%%%%%%%%%%%
% Cargando paquetes y definiendo opciones %
%%%%%%%%%%%%%%%%%%%%%%%%%%%%%%%%%%%%%%%%%%%
% Aquí puedes cargar los paquetes que vas a usar. La clase
% fime ya incluye babel, inputenc, graphicx y los de la AMS.
% Cargar un paquete está a tu libertad (y responsabilidad).
\usepackage{hyperref}
    \hypersetup{breaklinks=true,colorlinks=true,linkcolor=black,citecolor=black,urlcolor=black}
\usepackage[spanish,textsize=scriptsize]{todonotes}
\usepackage{multicol}

%%%%%%%%%%%%%%%%%%%%%
% Definiendo campos %
%%%%%%%%%%%%%%%%%%%%%
\def\titulo{Título en español}
\def\tituloEng{Title in english}
\def\autor{Nombre Apellidos}
\def\matricula{1234567}
\def\grado{Doctorado en\ldots}
\def\gradoEng{Doctor in \ldots}
% En caso de que el grado tenga orientación o especialidad llenar el siguiente
% campo dejando un ESPACIO INICIAL, en caso contrario, dejar vacío
\def\orientacion{}
% Coloca el mes con mayúscula inicial
\def\fecha{Mes 20XX}
\def\fechaEng{Month 20XX}

\def\asesor{Grado.\ Nombre Apellidos}
\def\asesorEng{Nombre Apellidos, Grado.}
\def\revisorA{Grado.\ Nombre Apellidos}
\def\revisorAEng{Nombre Apellidos, Grado.}
\def\revisorB{Dra.\ Nombre Apellidos}
\def\revisorBEng{Nombre Apellidos, Grado.}
% En el caso de que tu tesis sea de doctorado activa la variable cambiándola a \doctoradotrue
% y define tus otros dos revisores
\newif\ifdoctorado\doctoradotrue
\def\revisorC{Grado.\ Nombre Apellidos}
\def\revisorCEng{Nombre Apellido, Grado.}
\def\revisorD{Grado.\ Nombre Apellido}
\def\revisorDEng{Nombre Apellido, Grado.}
% El visto bueno siempre va
\def\vobo{Grado. Nombre Apellido}
\def\voboEng{Nombre Apellido, Grado.}

%%%%%%%%%%%%%%%%%%%%%%%
% Inicia el documento %
%%%%%%%%%%%%%%%%%%%%%%%
\begin{document}

\frontmatter
\pagestyle{main}

%%%%%%%%%%%%%%%%%%%%%%%%
% Primer portada: UANL %
%%%%%%%%%%%%%%%%%%%%%%%%
\thispagestyle{empty}

\begin{scshape}
\begin{center}
	{\Large \uanl} \\[5mm]
	{\large \fime} \\[5mm]
	{\large \depg}
	\vskip15mm
	\includegraphics[height=55mm]{uanl}
	\vskip12mm
	\begin{tabular}{p{12cm}}
		\centering
		{\large \titulo}
	\end{tabular}
	\vskip7mm
	{por}\\[7mm]
	{\large \autor}\\[7mm]
        % Si en tu posgrado la tesis no es opcional (como sí lo es en licenciatura), no modifiques esta línea:
	{como requisito parcial para obtener el grado de}\\[3mm]
	\MakeUppercase{\grado}\\
	\orientacion
	\vfill
	\fecha
\end{center}
\end{scshape}
 
%%%%%%%%%%%%%%%%%%%%%%%%%
% Segunda portada: FIME %
%%%%%%%%%%%%%%%%%%%%%%%%%
\newpage
\thispagestyle{empty}

\begin{scshape}
\begin{center}
	{\Large\uanl} \\[5mm]
	{\large\fime} \\[5mm]
	{\large\depg}
	\vskip16mm
	\includegraphics[height=50mm]{fime}
	\vskip16mm
	\begin{tabular}{p{12cm}}
		\centering
		{\large \titulo}
	\end{tabular}
	\vskip7mm
	{por}\\[7mm]
	{\large \autor}\\[7mm]
        % Si en tu posgrado la tesis no es opcional (como sí lo es en licenciatura), no modifiques esta línea:
	{como requisito parcial para obtener el grado de}\\[3mm]
	\MakeUppercase{\grado}\\
	\orientacion
	\vfill
	\fecha
\end{center}
\end{scshape}

%%%%%%%%%%%%%%%%%%%%%%%%%%%%%
% Carta del comité de tesis %
%%%%%%%%%%%%%%%%%%%%%%%%%%%%%
\newpage
\thispagestyle{empty}
\enlargethispage{5mm}

{\renewcommand{\baselinestretch}{1.1}\selectfont
\begin{center}\vspace*{-25mm}\hspace*{-10mm}
\begin{minipage}{170.5mm}
\hspace{-1.5mm}\includegraphics[height=20mm]{uanl}\hfill\raise1mm\hbox{\includegraphics[height=18.5mm]{fime}}
\hrule\vspace{0.5mm}
\scalebox{.5}{\MakeUppercase{\uanl}}\hfill\scalebox{.5}{\MakeUppercase{\fime}}\medskip
\end{minipage}
\vskip4mm{\sc\large\uanl\\\fime\\[3pt]\depg}\vskip6mm
\end{center}

Los miembros del Comité de Tesis recomendamos que la Tesis \textit{\titulo}, realizada por el alumno \autor, con número de matrícula \matricula, sea aceptada para su defensa como requisito parcial para obtener el grado de \grado\orientacion.
%\ifdoctorado\vskip10mm\else\vskip8mm\fi

\begin{center}
El Comité de Tesis\\
\ifdoctorado\vskip15mm\else\vskip25mm\fi

\ifdoctorado{%%%
\begin{tabular}{p{35mm}p{20mm}p{12mm}p{20mm}p{35mm}}
	\cline{1-2} \cline{4-5} \\ [-3mm]
	\multicolumn{2}{c}{\asesor} & & \multicolumn{2}{c}{\revisorA} \\
	\multicolumn{2}{c}{Asesora}  & & \multicolumn{2}{c}{Asesor} \\ [13mm]
	\cline{1-2} \cline{4-5} \\ [-3mm]
	\multicolumn{2}{c}{\revisorB} & & \multicolumn{2}{c}{\revisorC} \\
	\multicolumn{2}{c}{Revisora}   & & \multicolumn{2}{c}{Revisor} \\ [13mm]
	\cline{2-4} \\ [-3mm]
	& \multicolumn{3}{c}{\revisorD} & \\
	& \multicolumn{3}{c}{Revisora}   & \\ [17mm]
	& \multicolumn{3}{c}{Vo. Bo.} & \\ [13mm]
	\cline{2-4} \\ [-3mm]
	& \multicolumn{3}{c}{\vobo}   & \\
	& \multicolumn{3}{c}{Subdirector de Estudios de Posgrado}   & \\ &&&&
\end{tabular}
}\else{%%%
\begin{tabular}{p{37mm}p{21mm}p{12mm}p{21mm}p{37mm}}
	\cline{2-4}
	& \multicolumn{3}{c}{\asesor} & \\
	& \multicolumn{3}{c}{Asesor}  & \\[19mm]
	\cline{1-2} \cline{4-5}
	\multicolumn{2}{c}{\revisorA} & & \multicolumn{2}{c}{\revisorB} \\
	\multicolumn{2}{c}{Revisor}   & & \multicolumn{2}{c}{Revisor} \\
	& \multicolumn{3}{c}{Vo\. Bo.} & \\[100mm]
	\cline{2-4}
	& \multicolumn{3}{c}{\vobo}   & \\
	& \multicolumn{3}{c}{Subdirector de Estudios de Posgrado}   & \\ &&&&
\end{tabular}
}\fi%%%

\vfill

\snnl, \MakeLowercase{\fecha}

\end{center}}

% Dedicatoria

\thispagestyle{empty}
\vspace*{17mm}

\begin{flushright}
\begin{itshape}

Algo breve, como un minicuento.

\end{itshape}
\end{flushright}



\tableofcontents
\listoffigures
\listoftables

%Agradecimientos

\chapter{Agradecimientos}
\markboth{Agradecimientos}{}

A FIME, UANL, CONACYT y demás.

%Resumen

\chapter{Resumen}
\markboth{Resumen}{}

{\renewcommand{\baselinestretch}{1.1}\selectfont
{\setlength{\leftskip}{10mm}
\setlength{\parindent}{-10mm}

\autor.

Candidato para obtener el grado de \grado\orientacion.

\uanl.

\fime.

Título del estudio: \textsc{\titulo}.

\noindent Número de páginas: \pageref*{lastpage}.}

\paragraph{Introducción}
Lorem ipsum.

\paragraph{Antecedentes}
Lorem ipsum.

\paragraph{Hipótesis}
Lorem ipsum.

\paragraph{Preguntas de investigación}
\begin{itemize}
    \item ¿Pregunta 1?
    \item ¿Pregunta 2?
    \item ¿Pregunta 3?
\end{itemize}

\paragraph{Objetivo general}
Lorem ipsum.

\paragraph{Objetivos específicos}
\begin{itemize}
    \item Objetivo 1.
    \item Objetivo 2.
    \item Objetivo 3.
\end{itemize}

\paragraph{Método de estudio}
Lorem ipsum.

\paragraph{Contribuciones y conclusiones}
Lorem ipsum.

\paragraph{Trabajo a futuro}
Lorem ipsum.

\paragraph{Estructura de la tesis}
Lorem ipsum.

\vspace{15mm}

Firmas de asesores

\begin{tabular}{p{35mm}p{20mm}p{12mm}p{20mm}p{35mm}}
	\cline{1-2} \cline{4-5} \\ [-3mm]
	\multicolumn{2}{c}{\asesor} & & \multicolumn{2}{c}{\revisorA} \\
	\multicolumn{2}{c}{Coasesora}  & & \multicolumn{2}{c}{Coasesor} \\ [13mm]
\end{tabular}
%Resumen

\chapter{Abstract}
\markboth{Abstract}{}

{\renewcommand{\baselinestretch}{1.1}\selectfont
{\setlength{\leftskip}{10mm}
\setlength{\parindent}{-10mm}

\autor.

As a candidate to obtain the degree of \gradoEng\orientacion.

\uanl.

\fime.

Title of the study: \textsc{\tituloEng}.

\noindent Number of pages: \pageref*{lastpage}.}

\paragraph{Introduction}
Lorem ipsum

\paragraph{Background}
Lorem ipsum

\paragraph{Hipótesis}
Lorem ipsum

\paragraph{Research questions}
\begin{itemize}
    \item Question 1?
    \item Question 2?
    \item Question 3?
\end{itemize}

\paragraph{General objective}
Lorem ipsum.

\paragraph{Specific objectives}
\begin{itemize}
    \item Objective 1.
    \item Objective 2.
    \item Objective 3.
\end{itemize}

\paragraph{Methodology}
Lorem ipsum.

\paragraph{Contributions and conclusions}
Lorem ipsum.

\paragraph{Future work}
Lorem ipsum.

\paragraph{Thesis structure}
Lorem ipsum.

\vspace{15mm}

\begin{tabular}{p{35mm}p{20mm}p{12mm}p{20mm}p{35mm}}
	\cline{1-2} \cline{4-5} \\ [-3mm]
	\multicolumn{2}{c}{\asesorEng} & & \multicolumn{2}{c}{\revisorAEng} \\
	\multicolumn{2}{c}{Advisor}  & & \multicolumn{2}{c}{Advisor} \\ [13mm]
\end{tabular}

\mainmatter
\pagestyle{fime}

\chapter{Introduction}
\label{ch_Intro}

Lorem ipsum.
\chapter{Related literature}
\label{ch_Rel}

Lorem ipsum.
\chapter{Theoretical basis}
\label{ch_The}

Lorem ipsum.
\chapter{Methodology}
\label{ch_Met}

Lorem ipsum.
\chapter{Results}
\label{ch_Res}

Lorem ipsum.
\chapter{Conclusions}
\label{ch_Con}

Lorem ipsum.

\appendix
%%% Desbloquear si tienes apéndice
%\chapter{Este es un apéndice}

\section{Citas bibliográficas}

En princicpio tienes total libertad de incluir tu bibliografía con el entorno {\tt thebibliography} nativo de \LaTeX{} o mediante la herramienta \textsc{Bib}\TeX. En caso de que optes por esto último (recomendado), puedes usar alguno de los archivos {\tt mighelbib.bst} o {\tt mighelnat.bst} incluidos en el paquete {\tt Tesis-FIME}, pues sus diseños están basados en el estilo bibliográfico estándar del español, además de que armoniza con el estilo de tesis provisto por {\tt fime.cls}.

El estilo bibliográfico {\tt mighelbib} es numérico, es decir cita con un número entre corchetes, por ejemplo una cita \verb+\cite{Dan82}+ genera una etiqueta del tipo [13], mientras que el estilo {\tt mighelnat} es tipo autor-año y requiere que el paquete {\tt natbib} sea cargado (sin opciones) para su correcto funcionamiento, cita con el apellido del autor y el año, por ejemplo una cita \verb+\citet{Dan82}+ genera una etiqueta del tipo Dantzig (1982), mientras que una cita \verb+\citep{Dan82}+ genera una etiqueta del tipo (Dantzig, 1982).

Como muestra del estilo, unas citas: un libro clásico de programación lineal \citep{Baz04} y un documento histórico \citep{Dan82}. Para saber un poco más del uso de \textsc{Bib}\TeX, se puede leer \citep{Mat11}.

\section{Comillas}

El objetivo de esta sección era provocar otra página para que se vea el encabezado. Pero aprovechamos para decir que la clase {\tt fime.cls} carga el paquete {\tt babel} con la opción {\tt spanish}, por lo que cambiará automáticamente los dobles signos $<<$ y $>>$ por << y >>. Estas comillas angulares son las correctas en el idioma español, y son las que se usan en la clase {\tt fime.cls}, por lo que se sugiere sean las usadas en el texto cada que quieras <<entrecomillar>> algo.



\backmatter
\pagestyle{main}

% Aquí va la bibliografía, puedes usar el entorno de LaTeX (thebibliography)
% o la herramienta BibTeX. En caso de que optes por BibTeX, puedes usar
% alguno de los archivos de estilo (mighelbib.bst o mighelnat.bst) incluidos
% en el paquete, cuyos diseños armonizan con el diseño de tesis provisto por
% fime.cls. Para muestra, basta un botón:
\bibliographystyle{plainnat}
\bibliography{MiBiblio}

\label{lastpage}
%Autobiografia

\chapter*{Resumen autobiográfico}
\thispagestyle{empty}

\begin{center}
\autor

Candidato para obtener el grado de\\
\grado\\
\orientacion\bigskip

\uanl\\
\fime\bigskip

Tesis:\\
\textsc{\large\titulo}
\end{center}\bigskip

%Aquí va tu historia
Nací el X de mes de año en Ciudad, País; mis padres son Nombres Apellidos y Nombres Apellidos. En año egresé de la carrera de X en la Facultad de X de la Universidad X. En año me gradué como Maestro en X en la Facultad X de la Universidad X.

\end{document}

